\chapter{Thermodynamics essentials}\label{ch:thermo}

Heat is chapter~4's energy bookkeeping applied to $10^{23}$ actors:
the same conservation law, now tracking which share of the energy is
\emph{organised} (work) and which is \emph{scattered} (heat). Engines,
batteries and climate systems are all audits of that split.

% from-zero:begin
\begin{rulebox}[Before you start]
Chapters~4 and~7. \emph{Recap:} temperature measures the average kinetic energy of molecules; heat is energy in transit.
\end{rulebox}
\begin{wins}[After this chapter you can]
\begin{itemize}
\item use the ideal-gas law and equipartition;
\item apply the first law with a consistent sign convention;
\item draw processes on a $PV$ diagram;
\item compute the efficiency limit of a heat engine.
\end{itemize}
\end{wins}
\begin{principle}[How to think about this]
Heat is the energy of jostling molecules, and temperature is its average. The first law is energy conservation with heat added. The second says that heat engines always waste some, and Carnot shows how little they must.
\end{principle}
% from-zero:end

\section{Temperature, equilibrium and the gas law}\label{sec:th-gas}

Temperature is the scale-free way of saying ``equally spread''; ideal
gas ties the macro bookkeeping to it: $PV=nRT$. Microscopically
(figure~\ref{fig:wall}) pressure is momentum flux: each wall hit
deposits $2mv_x$, and summing the rate gives
$P=\tfrac{1}{3}nm\overline{v^{2}}/V$ --- the same Newton's law as
chapter~2, averaged.

\begin{figure}[htbp]\centering
\begin{tikzpicture}[scale=0.8,font=\scriptsize]
\draw[very thick] (4.6,-1.6) -- (4.6,1.6);
\foreach \x/\y/\a in {1.2/0.9/20, 2.4/-0.5/-30, 3.4/1.0/-15}{
\draw[->,thick,accent] (\x,\y) -- (4.5,\y+0.4);}
\node[align=left] at (2.6,-2.2) {each collision: impulse $2mv_x$;\\rate $\times$ impulse $=$ pressure};
\end{tikzpicture}

\caption{Kinetic theory: pressure is the sum of wall collisions.}
\label{fig:wall}
\end{figure}

\section{Internal energy and equipartition}\label{sec:th-internal}

For a monatomic gas each particle's average KE is $\tfrac32k_BT$, so
$E_{\rm int}=\tfrac32nRT$: temperature \emph{is} the mean per-particle
kinetic energy, and the distribution of speeds is not a spike but
figure~\ref{fig:mb} --- the Maxwell--Boltzmann shape the code samples.

\begin{figure}[htbp]\centering
\begin{tikzpicture}
\begin{axis}[width=8.2cm,height=4.8cm,xlabel={$v$},ylabel={$f(v)$},
xmin=0,xmax=3,ymin=0,ymax=1.0,font=\small]
\addplot[very thick,accent,domain=0:3,samples=80] {x*exp(-x^2)};
\addplot[dashed] coordinates {(1,0) (1,0.61)};
\node[align=center,font=\scriptsize] at (axis cs:2.0,0.55) {peak at $\sigma$;\\long tail $=$ fast few};
\end{axis}
\end{tikzpicture}

\caption{Maxwell--Boltzmann: most particles near $\sigma$, a fast
tail doing the chemistry.}
\label{fig:mb}
\end{figure}

\section{The first law: the fixed sign convention}\label{sec:th-first}

\begin{equation}
  \Delta E_{\rm int}=Q-W_{\rm by},\label{eq:first}
\end{equation}
with $Q$ the heat \emph{into} the gas and $W_{\rm by}$ the work done
\emph{by} it --- one convention, used throughout (a compression then
has $W_{\rm by}<0$: the gas gains energy).

\section{Processes on the $PV$ plane}\label{sec:th-proc}

Four archetypes (figure~\ref{fig:pv}): isochore ($W=0$), isobar
($W=P\Delta V$), isotherm ($PV$ const, $Q=W$), adiabat ($Q=0$,
$PV^{\gamma}$ const, steeper than the isotherm). Work is the area under
the process curve --- so a cycle's \emph{enclosed area} is its net work
(figure~\ref{fig:cycle}).

\begin{figure}[htbp]\centering
\begin{tikzpicture}
\begin{axis}[width=8.4cm,height=5.2cm,xlabel={$V$},ylabel={$P$},
xmin=0.5,xmax=4.2,ymin=0,ymax=2.6,font=\small,xtick=\empty,ytick=\empty]
\addplot[very thick,accent,domain=0.5:4,samples=50] {1/x};
\node[align=center,font=\scriptsize] at (axis cs:1.6,1.9) {isotherm\\$PV=\mathrm{const}$};
\addplot[very thick,amber,domain=0.5:2.2,samples=50] {2.2/x^1.667};
\node[align=center,font=\scriptsize,amber] at (axis cs:0.85,2.1) {adiabat};
\draw[dashed] (axis cs:1.4,0) -- (axis cs:1.4,2.5);
\node[align=center,font=\scriptsize] at (axis cs:1.5,2.35) {isochore};
\draw[dashed] (axis cs:0.5,0.55) -- (axis cs:4.0,0.55);
\node[align=center,font=\scriptsize] at (axis cs:3.4,0.75) {isobar};
\end{axis}
\end{tikzpicture}

\caption{The four process archetypes: the adiabat always beats the
isotherm's steepness.}
\label{fig:pv}
\end{figure}

\begin{figure}[htbp]\centering
\begin{tikzpicture}
\begin{axis}[width=8cm,height=4.8cm,xlabel={$V$},ylabel={$P$},
xmin=0.8,xmax=3.4,ymin=0.4,ymax=2.4,font=\small,xtick=\empty,ytick=\empty]
\addplot[very thick,accent,domain=1:3,samples=40] {1.4/x} ;
\addplot[very thick,amber,domain=1:3,samples=40] {2.0/x} ;
\draw[very thick] (axis cs:1,1.4) -- (axis cs:1,2.0);
\draw[very thick] (axis cs:3,0.467) -- (axis cs:3,0.667);
\addplot[fill=accent!15,domain=1:3,samples=40] {2.0/x} \closedcycle;
\node[align=center,font=\scriptsize] at (axis cs:2,1.15) {enclosed area\\$=$ net work};
\end{axis}
\end{tikzpicture}

\caption{Any clockwise cycle nets the enclosed area as work.}
\label{fig:cycle}
\end{figure}

\begin{example}[Isothermal expansion + audit]\label{ex:iso}
$2$\,mol of gas expands isothermally at $300$\,K from $1$ to $2$\,m$^{3}$.
Find $W$, and audit $Q$ and $\Delta E_{\rm int}$.
\par\smallskip
\textbf{Strategy.} $W=\int P\,dV=nRT\ln(V_2/V_1)$.
\begin{solution}
$W=2(8.314)(300)\ln2=3458$\,J. Ideal-gas $E_{\rm int}$ depends only on
$T$: $\Delta E_{\rm int}=0$, so the first law forces
$Q=W=3458$\,J --- every joule in becomes a joule out.
\checkmark
\end{solution}
\end{example}

\begin{example}[Two-process cycle: net work from the area]\label{ex:cyc}
A cycle: isobaric expansion $1\to2$\,m$^{3}$ at $2\times10^{5}$\,Pa,
then isochoric cooling back to the start state, then isobaric
compression $2\to1$\,m$^{3}$ at $1\times10^{5}$\,Pa, then isochoric
heating home. Net work per cycle?
\par\smallskip
\textbf{Strategy.} Enclosed rectangle on the $PV$ plane.
\begin{solution}
$W_{\rm net}=(2-1)\times10^{5}\times(2-1)=1.0\times10^{5}$\,J --- the
rectangle's area, no calculus needed. \checkmark
\end{solution}
\end{example}

\begin{example}[Adiabatic compression work]\label{ex:adiab}
$1$\,mol of air ($\gamma=1.4$, $C_V=\tfrac52R$) at $300$\,K, $0.02$\,m$^{3}$
is compressed adiabatically to $0.01$\,m$^{3}$. Find $T_2$ and the work
done \emph{on} the gas.
\par\smallskip
\textbf{Strategy.} $TV^{\gamma-1}$ const; $W_{\rm on}=-\Delta E_{\rm int}$.
\begin{solution}
$T_2=300(2)^{0.4}=396$\,K; $\Delta E_{\rm int}=nC_V\Delta T
=2.5(8.314)(96)=1995$\,J, all arriving as work done on the gas since
$Q=0$. \checkmark
\end{solution}
\end{example}

\section{Heat engines and the Carnot bound}\label{sec:th-carnot}

An engine (figure~\ref{fig:engine}) banks $Q_H$, pays $W$, dumps $Q_C$:
$\eta=W/Q_H=1-Q_C/Q_H$. The ideal-gas Carnot cycle (two isotherms + two
adiabats) computes $Q_C/Q_H=T_C/T_H$, giving the universal bound
\begin{equation}
  \eta\le1-\frac{T_C}{T_H},\label{eq:carnot}
\end{equation}
a rectangle on the $T$--$S$ plane (figure~\ref{fig:ts}) --- no real
engine does better.

\begin{figure}[htbp]\centering
\begin{tikzpicture}[scale=0.85,font=\scriptsize]
\draw[very thick,fill=amber!25] (0,1.6) ellipse (1.0 and 0.45);
\node at (0,1.6) {$T_H$};
\draw[very thick,fill=accent!15] (0,-1.6) ellipse (1.0 and 0.45);
\node at (0,-1.6) {$T_C$};
\draw[very thick,fill=white] (0,0) circle (0.55);
\node at (0,0) {eng.};
\draw[->,thick] (-0.3,1.15) -- (-0.3,0.6) node[midway,left]{$Q_H$};
\draw[->,thick] (-0.3,-0.6) -- (-0.3,-1.15) node[midway,left]{$Q_C$};
\draw[->,thick] (0.55,0) -- (1.3,0) node[right]{$W$};
\node[align=left] at (2.6,-0.9) {$\eta=W/Q_H$\\$\le 1-T_C/T_H$};
\end{tikzpicture}

\caption{The engine schematic: efficiency is the take vs the bank.}
\label{fig:engine}
\end{figure}

\begin{figure}[htbp]\centering
\begin{tikzpicture}
\begin{axis}[width=8cm,height=4.6cm,xlabel={$S$},ylabel={$T$},
xmin=0.5,xmax=3.5,ymin=0.5,ymax=3,font=\small,xtick=\empty,ytick=\empty]
\draw[very thick,accent] (axis cs:1,1) rectangle (axis cs:3,2.4);
\draw[->,thick] (axis cs:1.6,2.4) -- (axis cs:2.6,2.4);
\draw[->,thick] (axis cs:1.6,1.0) -- (axis cs:1.0,1.0);
\node[align=center,font=\scriptsize] at (axis cs:2,1.7) {Carnot:\\rectangle on $T$--$S$};
\end{axis}
\end{tikzpicture}

\caption{Carnot on $T$--$S$: a rectangle whose width is $\Delta S$.}
\label{fig:ts}
\end{figure}

\begin{example}[Carnot bound number]\label{ex:carnot}
An engine runs between $600$\,K and $300$\,K. Bound on $\eta$; if it
takes $1000$\,J per cycle, best possible $W$?
\par\smallskip
\textbf{Strategy.} Apply \eqref{eq:carnot}.
\begin{solution}
$\eta\le1-300/600=50\%$; $W\le500$\,J --- half the heat must leave as
$Q_C$ no matter the engineering. \checkmark
\end{solution}
\end{example}

\begin{example}[Isothermal entropy of expansion]\label{ex:entropy}
For \ref{ex:iso}'s expansion, find $\Delta S$ of the gas.
\par\smallskip
\textbf{Strategy.} $\Delta S=Q_{\rm rev}/T$.
\begin{solution}
$\Delta S=3458/300=11.5$\,J/K $=nR\ln2$ exactly --- the gas's states
multiplied by $2^{n}$ (each particle's volume doubled, in counting
terms). \checkmark
\end{solution}
\end{example}

\begin{lstlisting}[language=Python,caption={\texttt{code/ch17\_maxwell\_boltzmann.py} --- Kac-model thermalisation vs 2-D MB theory.},label={lst:mb}]
rx, ry = c*wx - s*wy, s*wx + c*wy   # rotate relative velocity
\end{lstlisting}

Output of a real run:
\begin{lstlisting}[numbers=none]
mean <v^2> = 1.0 (1.0: collisions conserve KE)
rms speed   = 1.0
mean speed  = 0.8876  2D MB theory = 0.8862
hist (0.4 bins) = [590, 1307, 1138, 657, 229, 69, 10, 0, 0, 0]
P(v>2.5) sim = 0.001  theory exp(-v^2/(2*0.5)) = 0.0019
\end{lstlisting}

Starting from \emph{every} particle at the same speed, random-angle
collisions alone build the MB shape: mean speed within $0.2\%$ of
theory and the tail in the right place.

\begin{example}[Molar heats: the gas audits its own degrees of freedom]\label{ex:molar}
Heat $2$\,mol of nitrogen gas through $\Delta T=100$\,K --- once at
constant volume, once at constant pressure. Find both heats and the
$\gamma$ they imply.
\par\smallskip
\textbf{Strategy.} Equipartition prices each quadratic degree at
$\tfrac12R$: diatomic $C_V=\tfrac52R$, and $C_P=C_V+R$ always.
\begin{solution}
$Q_V=nC_V\Delta T=2\cdot\tfrac52(8.314)(100)=4.16$\,kJ;
$Q_P=nC_P\Delta T=2\cdot\tfrac72(8.314)(100)=5.82$\,kJ;
their ratio $\gamma=C_P/C_V=7/5=1.40$.
The extra $1.66$\,kJ at constant pressure is exactly the expansion
work $nR\Delta T$ --- $C_P-C_V=R$ is the first law wearing the
molar-capacity label. Monatomic argon would run $3/2$ and $5/2$
($\gamma=1.67$): the capacities \emph{count} the molecule's degrees
of freedom, so a measured $\gamma$ is a microscopy result.
\checkmark
\end{solution}
\end{example}

\section{In engineering practice: the heat budget}\label{sec:th-eng}

Power plants are figure~\ref{fig:engine} at gigawatt scale (cooling
towers are the $T_C$ dump); refrigeration is the same loop reversed;
the Carnot bound is why no one sells a $90\%$ engine --- it isn't
engineering, it's counting.

\section{Common mistakes}\label{sec:th-mistakes}

\begin{tripwire}[$W$ sign flipped]
Our convention: $W_{\rm by}$ the gas. Compression is negative by
definition --- flip once and every subsequent sign follows.
\end{tripwire}
\begin{tripwire}[$Q=0$ read as $\Delta T=0$]
Adiabatic $\ne$ isothermal: no heat exchanged, yet work still changes
$T$ ($PV^\gamma$ steeper for a reason).
\end{tripwire}
\begin{tripwire}[Real efficiency vs Carnot conflated]
$\eta=1-Q_C/Q_H$ is bookkeeping that always holds; the bound
$T_C/T_H$ holds only for reversible cycles.
\end{tripwire}
\begin{tripwire}[Celsius in a ratio]
Every $T$ in \eqref{eq:carnot} and $PV=nRT$ is kelvin --- ratios in
${}^{\circ}$C are meaningless.
\end{tripwire}
\begin{tripwire}[Entropy mysticism]
$\Delta S=Q_{\rm rev}/T$ is counting states, not ``disorder'' ---
\ref{ex:entropy} computed it exactly.
\end{tripwire}

\section{Summary}\label{sec:th-summary}

\begin{table}[ht]\centering
\begin{tabularx}{\linewidth}{@{}lY@{}}
\toprule
\tablehead\thead{Idea} & \thead{One-line version}\\
\midrule
gas law & $PV=nRT$; $P=\tfrac13nm\overline{v^2}/V$ micro \\
energy & $\tfrac32k_BT$ per particle, monatomic \\
first law & $\Delta E_{\rm int}=Q-W_{\rm by}$ \\
processes & work $=$ area; adiabat $PV^\gamma$ steepest \\
Carnot & $\eta\le1-T_C/T_H$, universal \\
entropy & $\Delta S=Q_{\rm rev}/T$, counting states \\
\bottomrule
\end{tabularx}
\caption{Thermo cheat sheet.}
\label{tab:thsummary}
\end{table}

\begin{wins}[Definition of done]
You can now: read any process off the $PV$ plane, apply the first law
with the fixed sign, compute isothermal/adiabatic work, bound any
engine, and run the thermalisation check.
\end{wins}

\chapternotes{%
The order follows OpenStax vol.~2 chs.~1--4 \cite{openstax1}; MIT
5.60 for the deeper development; Feynman~I \S39--45 the kinetic
picture \cite{feynman1}. The numbers are the listing's real output.}

\begin{rulebox}[{Watch instead}]
\begin{enumerate}
\item \emph{Lec 1: Thermodynamics \& Kinetics} --- MIT 5.60, MIT
OpenCourseWare, 46:46, \url{https://youtu.be/kLqduWF6GXE} ---
\S\ref{sec:th-gas}--\S\ref{sec:th-first} (accessed 2026-10-11).
\item \emph{Kinetic Theory of Gases Part 1} --- MIT OpenCourseWare,
78:39, \url{https://youtu.be/BhVyiU_dWps} --- \S\ref{sec:th-gas} and
\S\ref{sec:th-internal} at full length (accessed 2026-10-11).
\item \emph{Carnot cycle and Carnot engine} --- Khan Academy, 20:53,
\url{https://youtu.be/aAfBSJObd6Y} --- \S\ref{sec:th-carnot} worked
(accessed 2026-10-11).
\end{enumerate}
\end{rulebox}

\section*{Further reading}
\addcontentsline{toc}{section}{Further reading}

\begin{itemize}
\item OpenStax \emph{College Physics 2e}, chs.~13--15 --- gas laws,
first law, engines: \url{https://openstax.org} (verified 2026-10-11).
\item MIT OCW \emph{5.60 Thermodynamics \& Kinetics} --- the full
course: at \url{https://ocw.mit.edu}, course ``5.60'' (verified
2026-10-11).
\item R.~P. Feynman, \emph{The Feynman Lectures on Physics} I,
\S39--45 --- kinetic theory at its clearest (in print).
\item Khan Academy, \emph{Thermodynamics} --- drills matching
\S\ref{sec:th-proc}--\S\ref{sec:th-carnot}: verified 2026-10-11 at
\url{https://www.khanacademy.org} (science $\to$ physics $\to$
thermodynamics).
\item PhET, \emph{Gas Properties} interactive --- pressure, temperature
and the microscopic picture: at \url{https://phet.colorado.edu}
(verified 2026-10-11).
\end{itemize}

\section*{Additional exercises}
\addcontentsline{toc}{section}{Additional exercises}

\begin{exercise}[easy]
A gas at $300$\,K occupies $2$\,m$^{3}$ at $10^{5}$\,Pa: find $n$.
(AP-style)
\end{exercise}
\begin{exercise}[easy]
Read the signs: gas gains $500$\,J heat and does $200$\,J work --- what
is $\Delta E_{\rm int}$? (AP-C)
\end{exercise}
\begin{exercise}[easy]
Which process is horizontal on the $PV$ plane? Which is vertical?
(F=ma)
\end{exercise}
\begin{exercise}[easy]
Why is $C_P>C_V$? One sentence, first law reasoning. (AP-C)
\end{exercise}
\begin{exercise}[easy]
A monatomic gas at $400$\,K: mean KE per particle in eV?
(AP-style)
\end{exercise}
\begin{exercise}[medium]
Adiabatic compression of $1$\,mol air ($\gamma=1.4$) halves the volume:
find $\Delta T$. (AP-C FR-style)
\end{exercise}
\begin{exercise}[medium]
A cycle encloses a $1$\,m$^{3}\times10^{5}$\,Pa square: net work per
cycle? (F=ma)
\end{exercise}
\begin{exercise}[medium]
Best efficiency between $800$\,K and room temperature? (F=ma)
\end{exercise}
\begin{exercise}[medium]
Entropy change of $1$\,mol expanding isothermally at $300$\,K to twice
the volume. (USAPhO-lite)
\end{exercise}
\begin{exercise}[medium]
Show $PV^\gamma=\mathrm{const}$ from the first law for an adiabat.
(derive)
\end{exercise}
\begin{exercise}[hard]
Free expansion ($Q=W=0$) still raises $S$: argue why without
violating the first law. (USAPhO-style)
\end{exercise}
\begin{exercise}[hard]
An engine between two \emph{finite} reservoirs that equilibrate:
find the maximum total work. (APhO-style)
\end{exercise}
\begin{exercise}[medium]
Sketch a refrigerator's $PV$ loop and label $Q_C$, $Q_H$, $W$.
(AP-C)
\end{exercise}
\begin{exercise}[medium]
Which is steeper on $PV$: adiabat or isotherm through the same point?
Prove it. (derive)
\end{exercise}
\begin{exercise}[medium]
Estimate the rms speed of a nitrogen molecule at room temperature.
(estimation)
\end{exercise}
\begin{exercise}[easy]
A monatomic gas's $C_V$ vs a diatomic's: which is larger and why does
rotation raise it? (AP-C)
\end{exercise}
\begin{exercise}[medium]
$2$\,mol of Ar is heated $50$\,K at constant volume, then again at
constant pressure: find both heats and the $\gamma$ they imply.
(F=ma)
\end{exercise}
\begin{exercise}[medium]
A heat pump keeps a house at $20^{\circ}$\,C against $-5^{\circ}$\,C
outside: bound the COP and estimate the electrical cost of $10$\,kW
of delivered heat. (estimation)
\end{exercise}
\begin{exercise}[hard]
Show the adiabatic TV-relation $TV^{\gamma-1}=$ const follows from
$PV^{\gamma}=$ const, then use it to find the temperature after
$10\times$ compression of air. (derive)
\end{exercise}
